\scanpage{15}
\uline{Exercice}\qquad Réalisation régulière d'un polygône.

Soit $\deleted{\unsure{V}}E$ plan affine euclidien sur $\mathbb{R}$ (avec metrique déf pos. sur l'espace $V$ des translations). Une réalisation rectiligne d'un polygone $(S,A,R)$ est dite \uline{régulière} si tous les cotes du polygône ont $=$ longueur\tikz[baseline=-.3ex,x=1em,y=1em]{
  \node[inner sep=0pt,font=\small,align=center] at (5.25,.05) {\inserted{\shortstack{et tous les angles\\itou}}};
  % The comma lies in the open left mouth of the bubble, beside longueur.
  \node[inner sep=0pt] at (.55,-.4) {,};
  % Two visible boundary strokes; do not complete the right edge beyond the scan.
  \draw[line width=.45pt] (.05,-.1) .. controls (.15,.45) and (-.9,.75) .. (-.65,1.35)
    .. controls (-.15,1.85) and (1.2,1.65) .. (2.2,1.55)
    .. controls (4.2,1.3) and (7.55,1.55) .. (9.85,1.35)
    .. controls (10.1,1.3) and (10.2,1.25) .. (10.3,1.1);
  \draw[line width=.45pt] (-.35,-.8) .. controls (.1,-.05) and (.8,.45) .. (1.15,-.1)
    .. controls (1.7,-.65) and (2.65,-.7) .. (3.5,-.8)
    .. controls (4.2,-1.3) and (6.8,-1.7) .. (7.55,-1.25)
    .. controls (8.45,-.8) and (9.85,-.55) .. (10.35,.05);
} régulière \deleted{et} \uline{normalisée} si la dite longueur est 1.

a) Existence, grâce à équation cyclotomique $z^n=1$ dans $\mathbb{C}$.

b) Prouver que si $(V\deleted{\unsure{,}}\supset P\deleted{\unsure{,}}\supset\deleted{\unsure{A}}S)$ et $(V'\supset P'\supset S')$ sont deux $n$-polygônes ($n\geq3$) munis d'une réal. rectiligne \uline{rég.} (resp normalisée) il existe une \deleted{\unsure{iso|isom}} similitude (resp \unsure{isom}) \deleted{\uline{unique}} de $V$ sur $V'$, qui transforme $P$ et $P'$, $S$ en $S'$;
\tikz[baseline=(cancel.base),x=1em,y=1em]{
  \node[inner sep=0pt,anchor=base west] (cancel) at (0,0) {\deleted{si $V\unsure{;|,}P,S$}};
  % The cancellation has a separate hooked return at its left end.
  \draw[line width=.4pt] (-.55,.25) .. controls (-.35,-.02) and (.12,-.06) .. (.62,.3);
  \draw[line width=.4pt] (-.55,.25) -- (-.12,.35);
} plus précisément, pour tt isom (comb) de $(\deleted{\unsure{P}}\unsure{S},A,R)$ sur $(S',A',R')$, $\exists\,\unsure{\vert|1}$ similitude
\tikz[baseline=(deV.base),x=1em,y=1em]{
  \node[inner sep=0pt,anchor=base west] (deV) at (0,0) {de $\deleted{\unsure{E}}V$};
  \node[inner sep=0pt,anchor=base west,font=\small] at (-1,1.75) {\inserted{(isom)}};
  \draw[->,line width=.4pt] (-.5,-.25) .. controls (-.05,.65) and (1.8,1.05) .. (3.1,1.2);
} sur $V'$ qui l'induise. \deleted{Montrer que}

\deleted{\unsure{Ex|Cor}}\quad\uline{Cor} Notion d'enveloppe affine (ou vectorielle, grâce : centre de
\tikz[baseline=(grav.base),x=1em,y=1em]{
  \node[inner sep=0pt,anchor=base west] (grav) at (0,0) {gravité)};
  \node[inner sep=0pt,anchor=base west,font=\small] at (-1.45,1.55) {\inserted{de l'ens des sommets}};
  \node[inner sep=0pt,anchor=base west] at (5.8,0) {d'un};
  % The independent closing parenthesis remains beneath the curved lead.
  \draw[line width=.4pt] (3.7,-.5) .. controls (3.95,.9) and (4.5,1.18) .. (5.2,1.12)
    .. controls (6.3,1.1) and (7.6,1.02) .. (8.45,1.18);
  \draw[line width=.4pt] (8.45,1.48) -- (8.3,.96) -- (9.1,1.5);
} polygône combinatoire.

\deleted{\unsure{Si|Soit}} Orientation du polygône $\longleftrightarrow$ orientation de son enveloppe

% Two detached dark points below the right end of the last line: writing versus surface trace remains uncertain.
\par\vspace{2.5\baselineskip}\noindent\hfill
\unsure{\tikz[baseline=-.3ex]{
  \fill (0,0) ellipse (.055em and .035em);
  \fill (.14em,.24em) ellipse (.035em and .055em);
}}\hspace*{2em}\par
% archive sheet number: 14
